\Titular%
{divulgacion}%
{Poesía e física: \textit{De catro a catro}}%
{Sebastián Táboas Pazo}%
{Como a física e a arte, a poesía de creación, son caras dunha mesma moeda.}%

\begin{multicols}{2}

\Cita[J. L. Borges]{La poesía es una necesidad, y he descubierto que no es rara. Además, he
descubierto que, muchas veces, caminando por la calle uno oye frases poéticas
dichas por personas que no saben que las frases son poéticas.}

\Cita[F. F. García Armesto]{Le preguntamos [a Rafael Dieste], pues: ¿por qué
siendo ante todo poeta (...) le han preocupado las matemáticas?}

\subsection*{Preámbulo. Ciencia e... arte?}

Ás veces, ao falarmos casualmente, atopamos esa postura dicotómica entre a
irreconciliabilidade das ciencias e das letras, unha guerra que moitas veces
levan a un discurso depreciativo entre o estudantado de diversas ramas. Sempre
sorprende aquel comentario onde na ciencia vese arte ou na arte, ciencia. Non
é, a caso, Kandinsky tan bo xeómetro como Hilbert?

A ciencia e a realidade tecnolóxica, que dela se deriva, sempre xoga un papel
fundamental nun traballo artístico, quizais só como marco, como contexto, mais
sempre hai quen fai da ciencia integramente arte. As vangardas de principios de
século XX presentaron unha revolución que respondía ao cambiante paradigma
tecno-científico, claro é revisando os ismos como o ultraísmo ou o futurismo.
Exemplo testemuñal é o manifesto \textit{Le Futurisme} \cite{futurismo} de
1909, <<\textit{IV- Declaramos que el esplendor del mundo se ha enriquecido de
una belleza nueva: la belleza de la velocidad. Un automóvil de carrera con su
vientre ornado de gruesas tuberías, parecidas a serpientes de aliento explosivo
y furioso...}>>.

% Un dos grandes poetas no século XX do noso país, creacionista excelso coma V.
% Huidobro, de educación x, chegouse a interesar moito polos traballos de
% Einstein (fixo ensaios), así o plasma no seu libro \textit{De catro a catro}

\subsection*{Manuel Antonio e Einstein}

Traiamos toda esta realidade artística aquí, a un Rianxo de comezos de século,
a un xove poeta galego, a Manuel Antonio.

Falamos agora dunha época de grandes cambios, marcada por esa filosofía da
sospeita\footnote{Paul Ricoeur tilda aos filósofos Marx (o traballo), Freud (o
subconsciente) e Nietzsche (os valores e a identidade individual) de ``mestres
da sospeita'' no seu libro \textit{Freud: unha interpretación da cultura} por
sinalar aquelas dimensións que subxacen baixo a identidade humana.}, auxe da
segunda revolución industrial, época convulsa, fin do imperialismo, vangardas
como resistencia e, na física, case completa (sentencia Kelvin), aparece un
recén doutorado, un xove Albert Einstein que publica (entre outros) un artigo
revolucionario, \textit{Zur Elektrodynamik bewegter Körper} (1905), até esas
famosas conferencias de 1915. Morría a física newtoniana, o tempo absoluto, e
nacía a moderna relatividade que atraía a todas esas cabezas famentas de novos
continentes, que atraía a todas esas vangardas.

\begin{figure}[H]
    \centering
    \includegraphics[width=\linewidth]{./revistas/005/imaxes/ma2_victor.jpg}
    \caption{Manuel Antonio, entusiasta do tabaco de pipa.}
    \label{im:manuelantonio}
\end{figure}

O rianxeiro mostrou sumo interese por todos estes traballos revolucionarios,
fervente lector das novidades científicas europeas (cónstanos que se volcara
moito na ciencia alemá \cite{MA_bib}), Nietzsche e os avances vangardistas,
entre estes, estaba moi ao tanto das novidades que trouxera Vicente Huidobro na
súa primeira visita a España en 1918 e da súa conferencia no Ateneo de Madrid
de 1921. Entre as súas lecturas científicas destaca a súa inmersión na
comprensión das investigacións e resultados de Einstein das que quedara
prendido cunha ollada. Todos estes detalles son coñecidos grazas á súa
correspondenza e á descrición da biblioteca do rianxeiro que fai García-Sabell.
Así, nunha carta que lle escribe a Cebreiro en 1922 figura:

\Cita{Merquei algún libro para me ``doutorar'' nela [a teoría da relatividade]
e agora ando tralo seguinte: Posto que a teoría da Relatividade demostra que o
Tempo e o Espazo e outras cousas son máis ou menos longas asegundo o sistema de
referenza, eu vou a desenrolar a miña autividade no medio en que o Tempo teña a
maisima longuedade e o Espazo a mínima.\cite{d4a4}}

Isto é sen lugar a dúbidas o xerme da xénes do poemario \textit{De catro a
catro}, unha xénese vencellada á física, aos recentes descubrimentos dos
científicos estranxeiros; poesía que se reduce á física. A marabillación do
poeta é tal que conclúe a súa carta cunha certa sorna adulación:

\Cita{Non se trata de parar o sol como Xosué, sinón de ir diante da cencia
alemá. ¿Como che quedou o corpo?\\ Este Dr.Einstein ``es un tigre'' como din os
``ches'', pero como non é galego ¡que se lle vai facer! ten que ficar sempre
por debaixo de nós.\cite{d4a4}}

Frecuentemente Manuel Antonio compartía e debatía ideas de toda índole, entre
elas literarias, con Dieste, acérrimo amigo de letras, tamén rianxeiro. Non foi
excepción desta idea xerme que figura na carta a Cebreiro, anos máis tarde
comentará:

\Cita{Siendo muy muchachos, me encontré con la sorpresa de que Manuel Antonio
lo suponía, sí [el espacio limitado, mínimo]. Me fue muy fácil persuadirle de
que ese límite, siendo una frontera, implicaba una continuación al otro lado y
de que aun suponiendo maciza e impenetrable esa continuación era imposible
concebirla sin espacio. Vi que acababa de abrirse como una nueva luz en su
consmovisión y que estaba por ello muy sorprendido y contentísimo.\cite{d4a4}}

O que sen dúbida foi fundamental para a creación de \textit{De catro a catro}
tal e como a coñecemos, coa súa escrita fragmentaria de espazo-temporalidade
tan marcada. Xa sexa dito de paso, segundo fun investigando e documentándome
para escreber este artigo que me roldaba a cabeza, atopeime con que Manuel
Antonio non era un exemplo tan xenial da física na poesía como si o é Rafael
Dieste coas matemáticas \cite{dieste_xeometria}.

\subsection*{Relatividade na teoría poética}

Baixo a pretensión de realizar unha breve análise da obra que encabeza este
artigo, afondaremos no ismo que máis a roza: \textit{o creacionismo}. Este
interese polas teorías de Einstein non era exclusiva do noso poeta, senón que
foi común en todos os creacionistas; xeneralizado a todas as vangardas e a
moitos ámbitos, moitas veces terxiversadas.

No escrito de 1926 \textit{El creacionismo} \cite{creacionismo}, Vicente
Huidobro sentencia que <<\textit{Lo realizado en la mecánica también se ha
hecho en la poesía}>>, sendo a nova teoría dinámica de Einstein un gran logro
da física, por que non conquistar tamén este continente na poesía? Fai,
principalmente, unha chamada a crear, que a única condición do poeta é
\textit{crear}, como xa declarara contundentemente na súa conferencia de xuño
de 1916 no Ateneo de Buenos Aires. Seguindo xa máis de cerca o criterio de
Gerardo Diego, <<\textit{o creador de imaxes comeza a crear por crear, non
describe, constrúe; non evoca, suxire}>> \cite{posibilidades}.

Podemos seguir un fío na pescuda de migallas relativistas nesta teoría poética.
Nese mesmo escrito de 1926 sinálase que <<\textit{La poesía creacionista
adquiere proporciones internacionales}>> en canto que a poesía para sublimar
non depende dunha lingua; é un invariante do sistema de referencia, non se
vislumbra aquí a caso unha teoría da relatividade case física?

A calidade máis notable é en realidade ese trato da dimensión espazo-temporal
no poema creacionista como se fose un tecido que depende do movemento, das
forzas, do observador, dunha clara implicación relativista que non se esforzan
en ocultar \cite{d4a4}. Explica Kathleen March \cite{march} que o xeito dos
poetas de utilizar estas nocións, ou máis ben, no caso de Manuel Antonio, é
<<\textit{Por medio do verbo}>>, explica que <<\textit{hai varias formas de
descompoñer a imaxe ou de construila para despois anulala. (...) os verbos, a
miúdo, pertencen a un campo de significado relacionado coa desaparición --ou á
destrución--. (...) Moitos deles apuntan a un tempo ou lugar alonxado, a unha
falta de precisión ou asequibilidade. Esto trae como consecuencia a
incapacidade de asir o obxecto poético que se atopa dotado dunha constante
dinamicidade.}>> A fin do emprego dos verbos é a anulación do tempo secuencial,
da cadea lóxica (como se se pretendese escribir fóra do cono de luz dende
dentro). Nun mesmo poema aparecen imaxes coespaciais cuxos verbos (que
descreben a acción) preséntanse tanto en pasado como en presente creando unha
organización allea ao tempo lóxico. <<\textit{Non hai cronoloxía obrigatoria,
aínda que o ollo do lector tende a desprazarse desde a parte superior da páxina
cara abaixo}.>> \cite{march}

\begin{figure}[H]
    \centering
    \includegraphics[width=\linewidth]{./revistas/005/imaxes/excelsior.jpg}
    \caption{Copia orixinal de \textit{Excelsior}, exposta na casa-museo.}
    \label{im:excelsior}
\end{figure}

Para combater esta verticalidade o poeta practicaba o caligrama, summun cubista
no poema, cuxo mellor exemplo é o poema \textit{Excelsior}. Formas combativas
máis lixeiras eran xogar coa separación e colocación dos versos, moi usuales
tamén aquí. Tamén é frecuente o xogo das voces, do cambio de persoa nos verbos
(ademais dos tempos), a aparición ``\textit{doutros}'' narradores para lograr
unha desubicación crono-espacial do lector; este fenómeno lembra a unha
multitude de sistemas de referencia para un mesmo suceso. Coa aparición de
tantas voces, en virtude de cal medimos o poema? cal é o noso sistema
absoluto?

March sinalaba o uso principal dos verbos para producir liñas de movemento e
tensión, porque se <<\textit{esixe o uso de termos concretos, ``visíbeis''}>>,
mais está claro que Manuel Antonio non se limitou só aos verbos. A ausencia,
espazos en branco, distancias aumentadas, tempos reducidos; voces, narradores
fragmentarios, multitude de sistemas de referencia isomorfos; recursos
sintácticos que soubo utilizar para \textit{crear} no absoluto sentido
creacionista, empapados da teoría da relatividade. Aventúrome a afirmar que,
dependendo do xogo sintáctico empregado, o poeta crea imaxes
\textit{space-like} ou imaxes \textit{time-like}.

Ao final o que máis aproveitan os creacionistas de Einstein é esa apertura do
sentido\footnote{O concepto da \textit{apertura do sentido}, polo menos así o
coñecín eu, aparece en \textit{Tensión e sentido} de M. Peyrou.} (a creación de
imaxes triplas, cuádruplas..., referirase Diego), utilizar a dimensionalidade
relativa (no seu sentido físico) para acadar, na creación de imaxes, unha
suxestión máxima. Inventar formas insólitas, mais sen rozar o surrealismo.
Aínda que moitas veces os literatos avogaran máis por unha construción
interpretativa da teoría da relatividade asíndoa a teorías do tempo máis
humanas, como as de Bergson (especialmente), Kant\footnote{Foi, en primeira
instancia, a concepción do espazo e do tempo de Kant a que puxo en xaque á
teoría do tempo e sistema de referencia absoluto de Newton e a que inspiraría
moi profundamente a Einstein nas súas investigacións.}, Prigogine (en tempos
modernos) e xunto con eses conceptos do devir e do eterno retorno de Nietzsche.

\subsection*{De catro a catro}

Rematamos agora o artigo comentando brevemente a obra do poeta rianxeiro, en
concreto o poemario de \textit{De catro a catro. Follas sin data dun diario
d'abordo}, que tanto se mencionou até agora. Manuel Antonio preséntanos un
mundo pechado no que non hai máis futuro nin máis pasado có da travesía. No
poema \textit{Travesía} xa se deixa clara as súas intencións
<<\textit{Troqueles reiterados/o reloxe e o Sol/alcuñaron moedas efímeras/que
repetían todas/a mesma cara e a mesma cruz}>>, pechándoo con

\Cita{O minuteiro\\\text{} \hspace{4em} (tic tac)\\ asumeu o compás das travesías.}

Esta concepción en certo xeito romántica (no sentido da Bildungsroman)
refórzase co poder da suxerencia que implicaba a mecánica relativista que ben
coñecía. Estamos perante un caso de simultaneidade \cite{march}; os tempos
están superpostos e atopamos que algo está a ocorrer neste mesmo instante aínda
que ocorrera na creación do universo. Non hai un antes nin un despois que non
estea vencellado á viaxe. É en virtude desta que as cousas acontecen; o sistema
de referencia é a viaxe en si mesma, ela é a medida das cousas; e a viaxe
carece dunha precisión cronolóxica, anulando a dimensión temporal. É evidente
tamén que <<non abonda coa lectura superficial da viaxe, vencellada ó
movemento, porque \textit{De catro a catro} é xa dende o título un poemario de
estatismo>>\footnote{Non atopei a referencia directa, pero en \cite{d4a4} é
atribuída a Elvira Souto.}, de aí que aquelas moedas repitan en bucle o mesmo
resultado.

Acábanse os ocos e non se pode un explaiar máis, tan só vomitar unhas frases
contundentes e crípticas. Este libro non só é física, a súa dimensión humana
resalta sobre todo en poemas como \textit{Sós} ou na monotonía do estatismo e
aínda así é imposible negar que nel habita un xerme. A comprensión do fondo (ou
ao que se refire Baudelaire con `\textit{fondo}') da obra de Manuel Antonio é
irrecoñecible sen a física do momento. Non podo seguir descifrando o poemario
nestas follas, non me queda máis que ir de catro en catro até que

\Cita{Encherémo-las velas\\ ca luz náufraga da madrugada}

\nocite{creacionismo}
\nocite{d4a4}
\nocite{dieste_xeometria}
\nocite{march}
\nocite{posibilidades}
\nocite{futurismo}

\printbibliography

\end{multicols}
